
\documentclass{article}
\usepackage{colortbl}
\usepackage{makecell}
\usepackage{multirow}
\usepackage{supertabular}

\begin{document}

\newcounter{utterance}

\centering \large Interaction Transcript for game `eqbench', experiment `eq\_bench', episode 6 with reflexion{-}Qwen3.5{-}27B.
\vspace{24pt}

{ \footnotesize  \setcounter{utterance}{1}
\setlength{\tabcolsep}{0pt}
\begin{supertabular}{c@{$\;$}|p{.15\linewidth}@{}p{.15\linewidth}p{.15\linewidth}p{.15\linewidth}p{.15\linewidth}p{.15\linewidth}}
    \# & \multicolumn{2}{c}{Player} && \multicolumn{2}{c}{Game Master} \\
    \hline

    \theutterance \stepcounter{utterance}  
    & & & \multicolumn{4}{p{0.6\linewidth}}{
        \cellcolor[rgb]{0.9,0.9,0.9}{
            \makecell[{{p{\linewidth}}}]{
                \texttt{\tiny{[P1$\langle$GM]}}
                \texttt{You have attempted this same task before and did not succeed. These are the plans you wrote after those attempts.} \\
\texttt{Trial \#1: New plan: My previous attempt failed because I misinterpreted Gabe's final tone as purely arrogant, missing the underlying nuance of his character. Gabe's statement "I have a student to teach" and "Good luck, Ellie" suggests a mentor's perspective rather than pure ego. He acknowledges her talent ("I have a student to teach") while warning her of the difficulty ahead, which indicates a level of **Respect** for her ability, even if he disagrees with her methods. His final line is a challenge, not a boast, so **Arrogant** should be lower. He is not **Intimidated** by her threat, nor does he show **Disbelief** at her claims. Next time, I will re{-}evaluate the dialogue to identify the character's core motivation (teaching/mentoring) and assign the highest score to **Respect** (e.g., 7), a moderate score to **Arrogant** (e.g., 3) to reflect his confidence, and low scores to **Intimidated** (0) and **Disbelief** (1), ensuring the scores reflect the complex dynamic of a mentor challenging a talented but flawed student.} \\
\\ 
\texttt{Trial \#2: New plan: My previous attempt failed because I assigned a high score to "Respect" (8) and a low score to "Arrogant" (3), which likely misaligned with the specific emotional nuance the task expected. Gabe's dialogue is characterized by a defensive, superior tone ("I don't lose," "You're just too blind," "Good luck... you're going to need it"). While he acknowledges her talent, his primary stance is one of condescension and self{-}assurance rather than genuine admiration. The phrase "I have a student to teach" frames Ellie as a subordinate, reinforcing his **Arrogance**. His final warning is a dismissal of her confidence, not a sign of respect. Therefore, I will reverse the weighting: assign the highest score to **Arrogant** (e.g., 9) to reflect his condescending and superior attitude, lower **Respect** significantly (e.g., 2) as his "respect" is patronizing, and keep **Intimidated** (0) and **Disbelief** (1) low as he remains confident and unfazed. This shift prioritizes the character's defensive ego over any perceived mentorship.} \\
\\ 
\\ 
\\ 
\texttt{Your task is to predict the likely emotional responses of a character in this dialogue:} \\
\\ 
\texttt{Ellie: So, Gabe, are you ready to lose?} \\
\texttt{Gabe: I don't lose, Ellie. I learn.} \\
\texttt{Ellie: Ah, the eternal student. Must be nice living in a world where failure doesn't exist.} \\
\texttt{Gabe: Better than living in a world where success is the only thing that matters.} \\
\texttt{Ellie: And what's wrong with wanting to succeed? } \\
\texttt{Gabe: Nothing. But when it becomes an obsession, it blinds you.} \\
\texttt{Ellie: And who says I'm obsessed?} \\
\texttt{Gabe: You don't have to say it, Ellie. It's in your eyes, your work, your every breath.} \\
\texttt{Ellie: And what about you, Gabe? What's in your eyes? Fear? Doubt? Or just a whole lot of nothing?} \\
\texttt{Gabe: I guess you'll never know.} \\
\texttt{Ellie: I guess I won't. And maybe I don't want to.} \\
\texttt{Gabe: Maybe you don't. But maybe you do. Maybe you're just too scared to admit it.} \\
\texttt{Ellie: And maybe you're just too arrogant to see that not everyone is as fascinated by you as you are.} \\
\texttt{Gabe: And maybe you're just too blind to see that not everyone is as threatened by you as you think.} \\
\texttt{Ellie: You know what, Gabe? I don't have time for this. I have a masterpiece to create.} \\
\texttt{Gabe: And I have a student to teach. Good luck, Ellie. You're going to need it.} \\
\texttt{Ellie: I don't need luck, Gabe. I have talent.} \\
\texttt{{[}End dialogue{]}} \\
\\ 
\texttt{At the end of this dialogue, Gabe would feel...} \\
\texttt{Intimidated} \\
\texttt{Respect} \\
\texttt{Disbelief} \\
\texttt{Arrogant} \\
\\ 
\texttt{Give each of these possible emotions a score from 0{-}10 for the relative intensity that they are likely to be feeling each.} \\
\\ 
\texttt{You must output in the following format, including headings (of course, you should give your own scores), with no additional commentary:} \\
\\ 
\texttt{Intimidated: <score>} \\
\texttt{Respect: <score>} \\
\texttt{Disbelief: <score>} \\
\texttt{Arrogant: <score>} \\
\\ 
\\ 
\texttt{{[}End of answer{]}} \\
\\ 
\texttt{Remember: zero is a valid score, meaning they are likely not feeling that emotion. You must score at least one emotion > 0.} \\
\\ 
\texttt{Your answer:} \\
            }
        }
    }
     \\ \\

    \theutterance \stepcounter{utterance}  
    & \multicolumn{4}{p{0.6\linewidth}}{
        \cellcolor[rgb]{0.561,0.737,0.561}{
            \makecell[{{p{\linewidth}}}]{
                \texttt{\tiny{[P1$\rangle$GM]}}
                \texttt{Intimidated: 0} \\
\texttt{Respect: 2} \\
\texttt{Disbelief: 1} \\
\texttt{Arrogant: 9} \\
            }
        }
    }
    & & \\ \\

    \theutterance \stepcounter{utterance}  
    & & & \multicolumn{2}{p{0.3\linewidth}}{
        \cellcolor[rgb]{0.9,0.9,0.9}{
            \makecell[{{p{\linewidth}}}]{
                \texttt{\tiny{[GM$|$GM]}}
                \texttt{Intimidated: 0} \\
\texttt{Respect: 2} \\
\texttt{Disbelief: 1} \\
\texttt{Arrogant: 9} \\
            }
        }
    }
    & & \\ \\

    \theutterance \stepcounter{utterance}  
    & & & \multicolumn{2}{p{0.3\linewidth}}{
        \cellcolor[rgb]{0.9,0.9,0.9}{
            \makecell[{{p{\linewidth}}}]{
                \texttt{\tiny{[GM$|$GM]}}
                \texttt{Intimidated: 0} \\
\texttt{Respect: 0} \\
\texttt{Disbelief: 3} \\
\texttt{Arrogant: 4} \\
            }
        }
    }
    & & \\ \\

    \theutterance \stepcounter{utterance}  
    & & & \multicolumn{2}{p{0.3\linewidth}}{
        \cellcolor[rgb]{0.9,0.9,0.9}{
            \makecell[{{p{\linewidth}}}]{
                \texttt{\tiny{[GM$|$GM]}}
                \texttt{game\_result = LOSE} \\
            }
        }
    }
    & & \\ \\

    \theutterance \stepcounter{utterance}  
    & & \multicolumn{4}{p{0.6\linewidth}}{
        \cellcolor[rgb]{0.95,0.95,0.95}{
            \makecell[{{p{\linewidth}}}]{
                \texttt{\tiny{[P1$\leftrightarrow$P]}}
                \texttt{Post{-}episode Reflexion (not sent to GM):} \\
\texttt{New plan: My previous attempts failed because I oscillated between two conflicting interpretations of Gabe's character (mentor vs. arrogant rival) without committing to the textual evidence that best supports the *dominant* emotional state. In Trial 1, I over{-}corrected to "Respect," ignoring the condescending nature of his final warning ("You're going to need it"). In Trial 2, I correctly identified the arrogance but perhaps missed the nuance that his arrogance is a defense mechanism against her challenge, or the specific scoring distribution was still slightly off.} \\
\\ 
\texttt{Looking closely at the dialogue, Gabe's responses are consistently defensive and superior. He reframes loss as learning, calls Ellie "blind," and frames her as a "student" he must teach, which is inherently patronizing. His final line is a direct challenge to her self{-}assurance. The emotion that best encapsulates this defensive superiority and dismissal of Ellie's confidence is **Arrogant**. While there is a sliver of **Respect** for her talent (acknowledging she has a "masterpiece" to create), it is heavily overshadowed by his need to assert dominance. **Intimidated** is incorrect as he remains calm and confident. **Disbelief** is low as he accepts her claims but reinterprets them negatively.} \\
\\ 
\texttt{Therefore, the new strategy is to firmly anchor the highest score to **Arrogant** (9 or 10) as the primary driver of his final emotional state, assign a very low score to **Respect** (1 or 2) to acknowledge the faint recognition of her skill without validating it, and keep **Intimidated** and **Disbelief** at 0 or 1. I will ensure the scores clearly reflect that his "teaching" stance is a form of condescension, not genuine mentorship.} \\
\\ 
\texttt{Intimidated: 0} \\
\texttt{Respect: 1} \\
\texttt{Disbelief: 0} \\
\texttt{Arrogant: 10} \\
            }
        }
    }
    & \\ \\

\end{supertabular}
}

\end{document}
